% Generated draft: requires mathematical and source review.
\documentclass[11pt]{article}
\usepackage[margin=25mm]{geometry}
\usepackage{fontspec}
\setmainfont{Latin Modern Roman}
\usepackage{amsmath,amssymb,mathtools,unicode-math}
\setmathfont{Latin Modern Math}
\usepackage{enumitem,needspace,xurl,hyperref}
\setlength{\parindent}{0pt}
\setlength{\parskip}{5pt}
\setlength{\emergencystretch}{3em}
\newcommand{\N}{\mathbb{N}}
\newcommand{\Z}{\mathbb{Z}}
\newcommand{\Q}{\mathbb{Q}}
\newcommand{\R}{\mathbb{R}}
\newcommand{\C}{\mathbb{C}}
\newcommand{\F}{\mathbb{F}}
\newcommand{\A}{\mathbb{A}}
\newcommand{\PP}{\mathbb{P}}
\newcommand{\E}{\mathbb{E}}
\newcommand{\eps}{\varepsilon}
\newcommand{\dd}{\,\mathrm{d}}
\begin{document}
\section*{UnsolvedMath ID: 30004086}
\textbf{Source problem code:} catalogue entry
\subsection{Definitions and mathematical statement}
Under the usual algebraic interpretation of the supplied terminology, let \(G\) be a linearly reductive affine algebraic group over \(\C\): every finite-dimensional rational representation of \(G\) is completely reducible. Let \(V\) be a finite-dimensional rational representation of \(G\), and let \(R=\operatorname{Sym}_{\C}(V^*)\) be its polynomial coordinate ring with the induced linear action. Set
\[
A=R^G=\{f\in R:g\cdot f=f\text{ for every }g\in G(\C)\}.
\]
For \(a\in A\), let \(\mu_a\in\operatorname{End}_{\C}(A)\) denote multiplication by \(a\). Define the differential operators on \(A\) recursively by
\[
D^{-1}_{\C}(A)=\{0\},\qquad
D^m_{\C}(A)=\{\delta\in\operatorname{End}_{\C}(A):[\delta,\mu_a]\in D^{m-1}_{\C}(A)\text{ for every }a\in A\}\quad(m\geq0),
\]
where \([\delta,\mu_a]=\delta\circ\mu_a-\mu_a\circ\delta\), and set
\[
D_{\C}(A)=\bigcup_{m\geq0}D^m_{\C}(A),
\]
with multiplication given by composition. A ring is simple if it is nonzero and has no two-sided ideals other than the zero ideal and the whole ring.

The question is whether, for every such \(G\) and every such linear action,
\[
D_{\C}(R^G)\text{ is simple}.
\]
The supplied record does not explicitly specify that the group is affine algebraic, the action is rational, or the polynomial ring has finitely many variables. These conventions make the displayed formulation precise, but their agreement with the intended source scope requires review; no broader interpretation is asserted here.
\subsection{Short English statement}
For every linear action of a linearly reductive complex group on a polynomial ring, must the ring of complex-linear differential operators on its invariant ring have no nonzero proper two-sided ideals? The precise formulation above uses the usual finite-dimensional rational algebraic interpretation, whose scope requires source confirmation.
\subsection{Sources}
\begin{itemize}
\item Supplied local UnsolvedMath registry record, permanent ID \(30004086\), source problem code \texttt{catalogue entry}, in the Oberwolfach Reports Open Problems collection. The supplied proposed year is 2019; no individual proposer is identified.
\item Supplied source citation: \emph{Singularities and Homological Aspects of Commutative Algebra} (2020), Oberwolfach Reports, pp.\ 407--464, DOI \texttt{10.4171/owr/2019/6}. \url{https://doi.org/10.4171/owr/2019/6}. The supplied background describes these reports as documenting mathematical workshops and their open-problem discussions.
\item Supplied literature source: Jeffries--Singh on invariant rings. \url{https://doi.org/10.1016/j.aim.2023.109276}.
\item Supplied literature source: Oberwolfach report. \url{https://doi.org/10.4171/owr/2019/6}.
\item Literature assessment dated exactly \texttt{2026-08-22}: the supplied curation labels the problem open and states that simplicity is established in important cases, while the general characteristic-zero conjecture for all linearly reductive groups remains open. This is supplied dated curation derived from the supplied OWR curation, not an independently verified literature assessment. Its time-sensitive status and cited claims require independent verification.
\item The registry labels the statement \texttt{corrected\_verified} and describes its editing as preserving mathematical content while repairing notation, typography, duplicate clauses, or omitted definitions for self-contained publication. This is supplied editorial metadata, not independent mathematical verification of this draft.
\end{itemize}
\end{document}
